\subsection{初等自同构}

设 \(\mathfrak{g}\) 是一个李代数。用 \(\operatorname{Aut}(\mathfrak{g})\) 记它的自同构群。若 \(x \in \mathfrak{g}\) 且 \(\operatorname{ad} x\) 幂零，则 \(e^{\operatorname{ad} x} \in \operatorname{Aut}(\mathfrak{g})\) （第一章，§6，第 8 节）。

定义 1. g 的形如 \(e^{ad x}\) （其中 ad x 幂零）的自同构的有限积称为 g 的初等自同构。g 的初等自同构群记作 \(\operatorname{Aut}_{e}(\mathfrak{g})\) 。

若 \(u \in \operatorname{Aut}(\mathfrak{g})\) ，则 \(ue^{\operatorname{ad}x}u^{-1} = e^{\operatorname{ad}u(x)}\) 。由此推出 \(\operatorname{Aut}_{e}(\mathfrak{g})\) 是 \(\operatorname{Aut}(\mathfrak{g})\) 的一个正规子群。若 k = R 或 C，则 \(\operatorname{Aut}_{e}(\mathfrak{g})\) 包含在 g 的内自同构群 \(\operatorname{Int}(\mathfrak{g})\) 中（第三章，§6，第 2 节，定义 2）。

* 在一般情形，\(\mathrm{Aut}_e(\mathfrak{g})\) 包含在代数群 \(\mathrm{Aut}(\mathfrak{g})\) 的单位元连通分支中。*

引理 1. 设 \(\mathrm{V}\) 是一个有限维向量空间，\(\mathfrak{n}\) 是 \(\mathfrak{a} = \mathfrak{gl}(\mathrm{V})\) 的一个由幂零元素组成的李子代数。

\begin{enumerate}
\def\labelenumi{(\roman{enumi})}
\item
  映射 \(x \mapsto \exp x\) 是从 \(\mathfrak{n}\) 到子群 \(N\) （\(\mathbf{GL}(V)\) 的，由幺幂元素组成）的双射（第二章，§6，第 1 节，注 4）。我们有 \(\mathfrak{n} = \log(\exp \mathfrak{n})\) 。映射 \(f \mapsto f \circ \log\) 是从 \(\mathfrak{n}\) 上的多项式函数代数到 \(\mathbf{N}\) 上的限制所成的代数（\(\operatorname{End}(\mathbf{V})\) 上的多项式函数的限制）的同构。
\item
  若 \(x \in \mathfrak{n}\) 且 \(a \in \mathfrak{a}\) ，则
\end{enumerate}

\[
(\exp \operatorname{ad} _ {\mathfrak {a}} x). a = (\exp x) a (\exp (- x)).
\]

\begin{enumerate}
\def\labelenumi{(\roman{enumi})}
\setcounter{enumi}{2}
\tightlist
\item
  设 \(V'\) 是一个有限维向量空间，\(n'\) 是 \(\mathfrak{gl}(V')\) 的一个由幂零元素组成的李子代数，\(\rho\) 是从 n 到 \(n'\) 的一个同态。设 \(\pi\) 是映射 \(\exp x \mapsto \exp \rho(x)\) （从 \(\exp n\) 到 \(\exp n'\) ）。那么 \(\pi\) 是一个群同态。
\end{enumerate}

由 Engel 定理，可以把 V 与 \(k^n\) 等同起来，使 \(\mathfrak{n}\) 是 \(\mathfrak{n}(n,k)\) （由对角线上为零的下三角矩阵组成的 \(\mathbf{M}_n(k)\) 的李子代数）的一个子代数。对 \(s \geq 0\) ，设 \(\mathfrak{n}_s(n,k)\) 是 \((x_{ij})_{1 \leq i,j \leq n} \in \mathbf{M}_n(k)\) 中使 \(x_{ij} = 0\) （对 \(i - j < s\) ）者所成之集。那么

\[
[ \mathfrak {n} _ {s} (n, k), \mathfrak {n} _ {s ^ {\prime}} (n, k) ] \subset \mathfrak {n} _ {s + s ^ {\prime}} (n, k)
\]

（第二章，§4，第 6 节，注），且 Hausdorff 级数定义了一个多项式映射 \((a,b)\mapsto\mathrm{H}(a,b)\) ，从 \(\mathfrak{n}(n,k)\times\mathfrak{n}(n,k)\) 到 \(\mathfrak{n}(n,k)\) （第二章，§6，第 5 节，注 3）；该映射使 \(\mathfrak{n}(n,k)\) 成为一个群（第二章，§6，第 5 节，命题 4）。由第二章 §6 第 1 节注 4，映射 \(x\mapsto\exp x\) （从 \(\mathfrak{n}(n,k)\) 到 \(1+\mathfrak{n}(n,k)\) ）与映射 \(y\mapsto\log y\) （从 \(1+\mathfrak{n}(n,k)\) 到 \(\mathfrak{n}(n,k)\) ）是互逆的双射且都是多项式的；由第二章 §6 第 5 节命题 3，若给 \(\mathfrak{n}(n,k)\) 赋予群律 \((a,b)\mapsto\mathrm{H}(a,b)\) ，并把 \(1+\mathfrak{n}(n,k)\) 视为 \(\mathbf{GL}_{n}(k)\) 的子群，则这些映射是群同构。于是引理的断言 (i) 和 (iii) 得出。设 \(x\in\mathfrak{n}\) 。用 \(\mathrm{L}_{x},\mathrm{R}_{x}\) 记映射 \(u\mapsto xu,u\mapsto ux\) （从 \(\mathfrak{a}\) 到 \(\mathfrak{a}\) ），它们交换且幂零。我们有 \(\mathrm{ad}_{\mathfrak{a}}x=\mathrm{L}_{x}-\mathrm{R}_{x}\) ，故对所有 \(a\in\mathfrak{a}\) ，

\[
\begin{array}{c} (\exp \operatorname{ad} _ {\mathfrak {a}} x) a = (\exp (\mathrm{L} _ {x} - \mathrm{R} _ {x})) a = (\exp \mathrm{L} _ {x}) (\exp \mathrm{R} _ {- x}) a \\ = \sum_ {i, j \geq 0} \frac {\mathrm{L} _ {x} ^ {i}}{i !} \frac {\mathrm{R} _ {- x} ^ {j}}{j !} a = (\exp x) a (\exp (- x)). \end{array}\tag{1}
\]

用引理 1 中的记号，\(\pi\) 称为 \(\exp \mathfrak{n}\) 的与给定表示 \(\rho\) （\(\mathfrak{n}\) 在 \(V'\) 上的）相容的线性表示。当 \(k\) 是 \(\mathbf{R}\) 、\(\mathbf{C}\) 或一个非离散完备超度量域时，由指数映射的性质有 \(\rho = L(\pi)\) （第三章，§4，第 4 节，命题 8 的推论 2）。

命题 1. 设 \(\mathfrak{g}\) 是一个李代数，\(\mathfrak{n}\) 是 \(\mathfrak{g}\) 的一个使 \(\operatorname{ad}_{\mathfrak{g}}x\) 对所有 \(x\in \mathfrak{n}\) 幂零的子代数。那么 \(e^{\operatorname{ad}_{\mathfrak{g}}\mathfrak{n}}\) 是 \(\operatorname{Aut}_e(\mathfrak{g})\) 的一个子群。

这由引理 1 (i) 立即得出。

特别地，若 \(\mathfrak{n}\) 是 \(\mathfrak{g}\) 的幂零根基，则 \(e^{\mathrm{ad}_{\mathfrak{g}}\mathfrak{n}}\) 是 \(\mathfrak{g}\) 的特殊自同构群（第一章，§6，第 8 节，定义 6）。

注. 1) 设 V 是一个有限维向量空间，g 是 \(\mathfrak{a} = \mathfrak{gl}(V)\) 的一个李子代数，x 是 g 的一个使 \(ad_{g}x\) 幂零的元素。那么存在 a 的一个幂零元素 n 使 \(ad_{a}n\) 延拓 \(ad_{g}x\) 。事实上，设 s, n 是 x 的半单分量与幂零分量；那么 \(ad_{a}s\) 和 \(ad_{a}n\) 是 \(ad_{a}x\) 的半单分量与幂零分量（第一章，§5，第 4 节，引理 2），故 \(ad_{a}s\) 和 \(ad_{a}n\) 保持 g 稳定，且 \(ad_{a}s|g\) 和 \(ad_{a}n|g\) 是 \(ad_{g}x\) 的半单分量与幂零分量；因此 \(ad_{g}x = ad_{a}n|g\) ，这就证明了我们的断言。注意到引理 1 (ii)，g 的每个初等自同构都延拓为 a 的一个形如 \(u \mapsto mum^{-1}\) （其中 \(m \in \mathbf{GL}(V)\) ）的自同构。

\begin{enumerate}
\def\labelenumi{\arabic{enumi})}
\setcounter{enumi}{1}
\tightlist
\item
  设 V 是一个有限维向量空间。对每个 \(g \in \mathbf{SL}(V)\) ，用 \(\varphi(g)\) 记自同构 \(x \mapsto gxg^{-1}\) （\(\mathfrak{gl}(V)\) 的）。那么
\end{enumerate}

\[
\operatorname{Aut} _ {e} (\mathfrak {g l} (V)) = \varphi (\mathbf {S L} (V)).
\]

事实上，由 (1)，\(\operatorname{Aut}_{e}(\mathfrak{gl}(V))\) 包含在 \(\varphi(\mathbf{SL}(V))\) 中，而反向包含由《代数》第三章 §8 第 9 节命题 17 及 (1) 得出。类似的论证表明 \(\operatorname{Aut}_{e}(\mathfrak{sl}(V))\) 是 \(\varphi(\mathbf{SL}(V))\) 的元素在 \(\mathfrak{sl}(V)\) 上的限制所成之集。

